% 相关工作
\chapter{相关工作}\label{chap:related-work}

本文横跨晶圆级集成、设计空间探索方法论与无阻塞网络理论三个领域，
相关工作也沿这三条线梳理，最后给出本文在三者交叉处的定位。

\section{晶圆级系统与晶圆级交换机}\label{sec:rw-wafer}

晶圆级集成的工业可行性已由若干标志性系统确立。
Tesla Dojo 的 Training Tile 以 InFO-SoW 封装集成了 25 个 D1 裸片，采用 2D Mesh + Z-Plane 拓扑~\cite{tesla2022dojo}；
Cerebras WSE 系列以单晶圆级芯片集成近百万处理核心，2D Mesh 提供数十 PB/s 的片上聚合带宽~\cite{cerebras2022wse2}。
这些系统确立了晶圆级集成作为规模化计算平台的可行性，但它们面向计算负载，而非网络交换。

晶圆级交换机作为独立研究方向由若干架构工作确立。
Chen 等~\cite{chen2024waferscale}首次系统量化了晶圆级交换机的 radix 提升潜力，
指出仅面积约束下可达传统芯片级交换机 32 倍的端口密度，并进一步识别出内部带宽、外部带宽与功率密度是真正的瓶颈；
Feng 和 Ma~\cite{feng2024switchless}提出「无交换机 Dragonfly」架构，以晶圆级高密度互联替代传统高 radix 物理交换机，
并给出死锁自由的最小/非最小路由；
Wan 等~\cite{wan2024architectural}在 2D Mesh 物理拓扑上叠加 Butterfly Fat-Tree 逻辑拓扑，实现 896 端口晶圆级交换系统；
Yang 等~\cite{yang2025pd}的 TickTock 框架面向 LLM 训练做物理-逻辑拓扑联合协同设计。
这些工作各自针对晶圆级交换的特定维度（radix 上限、特定拓扑的路由、逻辑-物理映射、计算负载协同设计）提供了深刻见解。

它们与本文的差别在于\textbf{分析方式}：上述工作多为自顶向下的上限估计或针对特定拓扑的架构设计；
本文则是自底向上的约束耦合判定——给定一个「拓扑 + 封装 + 散热」组合，判断它是否可行、瓶颈在哪、往哪改。
二者互补而非替代。

\section{DSE 方法论}\label{sec:rw-dse}

Chiplet 与晶圆级的设计空间探索工具近年快速涌现。
FPIA~\cite{jiao2024fpia}在可重用硅 interposer 上实现 field-programmable 互连 fabric 的 chiplet 级布局布线自动化，
其布线网格模型是本文布线约束的参考来源；
RapidChiplet~\cite{iff2025rapidchiplet}用分析模型替代 cycle-accurate 仿真做 chiplet 间互连的快速 DSE；
FireLink~\cite{li2025firelink}面向 chiplet 全栈 DSE（微架构 + 互连 + PPA + 成本），用决策树加速搜索；
MFIT~\cite{mfit2025}提供 2.5D/3D chiplet 系统的多保真热建模，是本文热 L0/L1 梯度配置的参考。
这些工具各自聚焦于单一物理维度（布局布线、互连性能、成本或热），
或面向小规模 chiplet 集成，缺乏跨层次（性能 + 几何 + 功耗 + 热）的联立判定能力。

与之相对，cycle-accurate 互连仿真器（如教科书~\cite{dally2004principles}所述的方法学）
单点评估需数分钟量级，对 $10^6$ 级设计空间的扫描不可行。
本文的 LP-based 方法在数学性质上与上述启发式/仿真方法有根本不同：
它不是启发式搜索，而是对每个设计点的可行性给出线性规划意义下的全局判定，
且借助事后闭式的影子价格（第~\ref{chap:sensitivity}章）提供精确的灵敏度信息——这是启发式方法无法提供的。

\section{无阻塞网络理论}\label{sec:rw-nonblocking}

无阻塞的概念源于经典交换理论：严格无阻塞（SNB）、宽感无阻塞（WSNB）与可重排无阻塞（RNB）
构成递进的松弛，其核心结果（如 Clos 三阶网络的 $m\ge 2n-1$ 与 $m\ge n$ 条件）是交换网络设计的经典基石，
系统性论述见 Dally 与 Towles 的教科书~\cite{dally2004principles}。
这一谱系关心的是\textbf{路径存在性}（能否连起来），不涉及本文的物理约束（bump、功耗、热）。

与本文性能侧最接近的是两个数学工具。
其一是 Birkhoff--von Neumann 定理~\cite{birkhoff1946tres}：双随机矩阵是置换矩阵的凸组合，
它把最坏流量模式归约到置换顶点，是本文排列模式集合 $\mathcal{R}$ 的数学来源。
其二是 Valiant 负载均衡路由~\cite{valiant1982scheme}：把任意流量先随机打散到中间节点再路由，
以两倍路径代价换取对抗性流量下的负载均衡，是本文「最优自适应路由压包络」的候选路径集来源之一。
Dragonfly 拓扑的「无阻塞」则是一个仿真驱动的经验概念——在均匀随机流量下 minimal routing 达到接近 100\% 的注入吞吐，
其平衡条件（$a=2p$、$h=p$）由 Kim 等~\cite{kim2008technology}给出，
这是比经典交换理论弱、但更贴近 HPC 实践的保证。

在方法学上最接近本文的是 Ngo 等的 LP 对偶框架~\cite{ngo2010lp-nonblocking}：
他们把无阻塞条件表示为 primal LP、取对偶、手构造对偶可行解，从而一次性推出各类交换网络的无阻塞充分条件。
但 Ngo 的方法处理的是\textbf{纯拓扑}无阻塞（Clos/Banyan 及其变体），「无阻塞」是二元的路径存在性；
本文则把统一的维度从「不同交换网络的非阻塞条件」扩展到「性能 + 几何 + 功耗 + 热」——
物理约束使「无阻塞」从二元判定变为「能撑到多高带宽」的连续量 $B^*$。

\section{本文的定位}\label{sec:rw-position}

综合上述三条线：晶圆级系统确立了平台、晶圆级交换机确立了问题、DSE 方法论与无阻塞理论各自提供了工具，
但\textbf{没有任何一项工作把性能、几何、功耗、热在同一组变量上联立为一个线性规划}。
本文填补的正是这一交叉空白：以负载包络 $\mathbf{L}$ 为枢纽，
把「最坏排列下的无阻塞潜能」与「按物理位置的约束」统一到单个可行性 LP，
并以 $B^*$ 与影子价格给出「可行/不可行」之外的瓶颈定位与投资优先级。
